\documentclass[exercice]{thedocument}%
\usepackage{texmaths-tikz}%
\startdatakeys
\source{SESAMATHS2016C4}
\cycle{4}
\competencesDuSocle{}
\automatismes{}
\objectifsApprentissage{}
\enddatakeys
\begin{document}%
    Les points \mathspoint{A}, \mathspoint{D} et \mathspoint{B} sont alignés.
    \par\smallskip\noindent
    En justifiant soigneusement votre réponse, calculer la mesure de l'angle
    $\widehat{\mathspoint{ABC}}$.
    \begin{center}
        \begin{tikzpicture}[scale=1.5,rotate=-15]
            \coordinate[label=below right:\mathspoint{A}](A) at (2,0);
            \coordinate[label=below left:\mathspoint{B}](B) at (-3,0);
            \coordinate[label=above:\mathspoint{C}](C) at (60:1.5);
            \coordinate[label=below left:\mathspoint{D}](D) at (0,0);
            \tkzFillAngle[\currentcolor!50, size=0.5cm](C,A,D)
            \tkzMarkAngle[size=0.5cm](C,A,D)
            \tkzLabelAngle[pos=0.75](C,A,D){$35^\circ$}
            \nextcolor
            \tkzFillAngle[\currentcolor!50, size=0.5cm](A,D,C)
            \tkzMarkAngle[fill=\currentcolor,size=0.5cm](A,D,C)
            \tkzLabelAngle[pos=0.75](A,D,C){$85^\circ$}
            \nextcolor
            \tkzFillAngle[\currentcolor!50, size=0.5cm](D,C,A)
            \tkzMarkAngle[fill=\currentcolor,mark=|,size=0.5cm](D,C,A)
            \tkzFillAngle[\currentcolor!50, size=0.5cm](B,C,D)
            \tkzMarkAngle[fill=\currentcolor,mark=|,size=0.5cm](B,C,D)
            \draw(A)--(B)--(C)--cycle;
            \draw(C)--(D);
        \end{tikzpicture}
    \end{center}
    \answerspace{4cm}
\end{document}